%%%%%%%% ICML 2026 LATEX SUBMISSION FILE %%%%%%%%%%%%%%%%%

\documentclass{article}

\usepackage{microtype}
\usepackage{graphicx}
\usepackage{subcaption}
\usepackage{booktabs}
\usepackage{hyperref}

\newcommand{\theHalgorithm}{\arabic{algorithm}}

\usepackage[accepted]{icml2026}

\usepackage{amsmath}
\usepackage{amssymb}
\usepackage{mathtools}
\usepackage{amsthm}
\usepackage{algorithm}
\usepackage{algorithmic}

\usepackage[capitalize,noabbrev]{cleveref}

%%%%%%%%%%%%%%%%%%%%%%%%%%%%%%%%
% THEOREMS
%%%%%%%%%%%%%%%%%%%%%%%%%%%%%%%%
\theoremstyle{plain}
\newtheorem{theorem}{Theorem}[section]
\newtheorem{proposition}[theorem]{Proposition}
\newtheorem{lemma}[theorem]{Lemma}
\newtheorem{corollary}[theorem]{Corollary}
\theoremstyle{definition}
\newtheorem{definition}[theorem]{Definition}
\newtheorem{assumption}[theorem]{Assumption}
\theoremstyle{remark}
\newtheorem{remark}[theorem]{Remark}

\icmltitlerunning{WB-CoMerge: Wasserstein-Barycenter Preconditioned Test-Time Model Merging}

\setlength{\textfloatsep}{9pt plus 2pt minus 2pt}
\setlength{\dbltextfloatsep}{9pt plus 2pt minus 2pt}
\setlength{\abovedisplayskip}{5pt plus 1pt minus 2pt}
\setlength{\belowdisplayskip}{5pt plus 1pt minus 2pt}
\setlength{\abovedisplayshortskip}{3pt}
\setlength{\belowdisplayshortskip}{3pt}

\begin{document}

\twocolumn[
  \icmltitle{WB-CoMerge: Wasserstein-Barycenter Preconditioned Co-acting\\
    Test-Time Model Merging for Non-Stationary Streams}

  \begin{icmlauthorlist}
    \icmlauthor{Anonymous Researcher}{equal,yyy}
  \end{icmlauthorlist}

  \icmlaffiliation{yyy}{Department of Computer Science, Autonomous University, Location, Country}
  \icmlcorrespondingauthor{Anonymous Researcher}{anonymous@edu.org}

  \icmlkeywords{Model Merging, Test-Time Adaptation, Optimal Transport, Wasserstein Barycenters}

  \vskip 0.3in
]

\printAffiliationsAndNotice{}

\begin{abstract}
  Test-time model merging (TTMM) has emerged as an efficient and data-free paradigm to specialize neural networks on non-stationary, unlabeled test streams. Existing state-of-the-art frameworks leverage Mixture-of-Gaussians (MoG) formulations to fuse the running statistics of Batch Normalization (BN) layers across diverse expert models. In this paper, we identify a fundamental flaw in MoG BN fusion: \emph{Variance Inflation}. When expert models specialize on highly dissimilar domains, their channel-wise mean activations differ significantly, causing the MoG variance to inflate proportionally to the squared distance between task means. Normalizing by this inflated variance excessively compresses the feature scales of both expert pathways, leading to representational collapse. To resolve this, we propose \textbf{WB-CoMerge} (Wasserstein-Barycenter Co-acting Model Merging), which models the channel-wise feature distributions as Gaussian probability measures and computes their geodesic interpolation under the Wasserstein-2 metric. By design, Wasserstein-Barycenter BN (WB-BN) completely eliminates the variance inflation term, preserving physical feature scales. We further extend this to \textbf{D-WB-CoMerge}, which dynamically blends static expert barycenters with streaming empirical batch statistics. Extensive evaluations on non-stationary vision streams demonstrate that WB-CoMerge prevents representational collapse, substantially outperforming existing SOTA baselines.
\end{abstract}

\section{Introduction}
\label{sec:intro}

Deep neural networks are typically deployed under the assumption that the test-time data distribution matches the training distribution. However, in real-world scenarios, models frequently encounter non-stationary, non-iid streams where the data distribution shifts dynamically and unpredictably~\cite{schneider2020improving}. Test-time adaptation (TTA) approaches seek to mitigate this domain shift on-the-fly using unlabeled test inputs~\cite{wang2020tent,liang2020shot}. While effective, standard TTA methods rely on backpropagation to update a large fraction of the model's parameters at inference time, which introduces substantial computational overhead, latency, and memory footprint, making deployment on resource-constrained edge devices highly challenging.

To amortize the cost of test-time adaptation, Test-Time Model Merging (TTMM) has emerged as a highly efficient alternative~\cite{hubotter2025local}. Instead of training or adapting deep architectures from scratch during inference, TTMM maintains a library of pre-trained "expert" neural networks (fine-tuned from a shared base initialization) and dynamically merges their weights in the parameter space on-the-fly. By routing the incoming unlabeled inputs to the most relevant experts and interpolating their parameters, TTMM can handle complex multi-task streams with a single forward pass, requiring zero backpropagation for basic routing and minimal gradient updates for local alignment.

A critical challenge in data-free model merging is the alignment of internal activations. When independent expert models are merged linearly, their internal feature spaces can drift, leading to severe activation mismatches and representational collapse~\cite{jordan2022repair}. To address this, state-of-the-art model merging frameworks, such as BK-CoMerge~\cite{anonymous}, utilize Mixture-of-Gaussians (MoG) moment matching to fuse Batch Normalization (BN) running mean and running variance buffers. 

In this work, we reveal a fundamental scientific and mathematical limitation in current MoG-based BN fusion strategies. Specifically, we show that when expert models specialize on highly distinct tasks (e.g., standard digit recognition vs.\ dense apparel classification), their internal features map to widely divergent channels and spatial centroids. Under an MoG assumption, the fused variance is calculated as:
\begin{equation}
  \sigma^2_{\text{MoG}} = \sum_k w_k \sigma_k^2 + \sum_k w_k (\mu_k - \mu_{\text{MoG}})^2
\end{equation}
where $\mu_k$ and $\sigma_k^2$ are the running mean and variance of expert $k$, and $w_k$ is the dynamic routing coefficient. Crucially, the second term $\sum_k w_k (\mu_k - \mu_{\text{MoG}})^2$ acts as a \emph{Variance Inflation} factor that scales quadratically with the distance between the experts' feature centroids $(\mu_1 - \mu_2)^2$. When the expert means are far apart, this inflation term dominates, causing the fused variance to explode. Consequently, when the merged model normalizes incoming activations by this artificially inflated standard deviation, the feature magnitudes of both pathways are excessively compressed (squashed), pushing them into the flat saturation regions of activation functions (such as ReLU). This leads to immediate representational collapse and catastrophic accuracy drops.

To overcome this limitation, we introduce \textbf{WB-CoMerge} (Wasserstein-Barycenter Co-acting Test-Time Model Merging). Instead of treating activations as a mixture of disjoint subsets, we treat the channel-wise feature distributions of each expert as independent Gaussian probability measures and find their geodesic interpolation under the Wasserstein-2 optimal transport metric~\cite{agueh2011barycenters,peyre2019computational}. Under the Wasserstein-2 metric, the barycenter of 1D Gaussian distributions (which perfectly models the channel-wise tracking of Batch Normalization) is a Gaussian distribution whose mean is the linear combination of means and whose standard deviation is the linear combination of standard deviations:
\begin{equation}
  \mu_{\text{WB}} = \sum_k w_k \mu_k, \quad \sigma_{\text{WB}} = \sum_k w_k \sigma_k
\end{equation}
By directly interpolating standard deviations rather than variances, Wasserstein-Barycenter BN (WB-BN) **completely eliminates the bimodal variance inflation term**. This mathematically elegant formulation preserves the physical scale of the internal feature maps, preventing representational collapse under any task-distance regime.

Furthermore, we extend our method to \textbf{D-WB-CoMerge} (Dynamic Wasserstein-Barycenter), which dynamically blends the static expert Wasserstein-barycenters with the empirical channel-wise mean and standard deviation of the incoming streaming batch on-the-fly. This allows the model to absorb non-stationary environmental corruptions (such as severe pixel-wise additive noise) in real-time, providing outstanding noise robustness.

Our key contributions are summarized as follows:
\begin{itemize}
  \item We mathematically expose the phenomenon of \emph{Variance Inflation} in existing Mixture-of-Gaussians Batch Normalization fusion methods, proving it is the primary driver of activation collapse when merging highly divergent experts.
  \item We propose \textbf{WB-CoMerge}, which utilizes Wasserstein-2 Barycenter interpolation of Batch Normalization running statistics to completely eliminate the variance inflation term.
  \item We design \textbf{D-WB-CoMerge}, a dynamic variant that fuses static expert barycenters with streaming empirical batch statistics to handle severe environmental corruptions.
  \item We conduct exhaustive evaluations on sequential multi-task vision streams, demonstrating that our proposed Wasserstein-barycenter methods significantly outperform existing SOTA test-time model merging baselines.
\end{itemize}

\section{Related Work}
\label{sec:related}

\textbf{Model Merging and Arithmetic.} Model merging aims to combine multiple task-specific neural networks into a single multi-task model without retraining~\cite{wortsman2022model,ainsworth2023git}. Standard averaging in weight space is highly effective when models share a pre-training initialization, as task-specific parameters remain in a shared loss basin~\cite{choshen2022fusing}. Beyond simple weight averaging, recent works introduce Fisher information weighting~\cite{matena2022merging}, task arithmetic subtraction and addition~\cite{ilharco2022editing}, and interference resolution via sign consensus~\cite{yadav2023ties}. However, these offline merging methods assume a static multi-task distribution and require either validation datasets or manual weight tuning, rendering them unsuitable for streaming, data-free test-time scenarios.

\textbf{Test-Time Adaptation (TTA).} Test-time adaptation specializes pre-trained models to non-stationary environments using unlabeled test batches on-the-fly~\cite{wang2020tent,schneider2020improving}. Traditional methods utilize entropy minimization~\cite{wang2020tent}, source-free domain adaptation~\cite{liang2020shot}, or single-sample self-supervision~\cite{zhang2022memo}. While effective, these methods rely on backpropagating gradients through the entire network (or submodules), which incurs severe latency and memory overhead on resource-constrained hardware. Our method bypasses heavy parameter training by dynamically merging pre-trained expert parameters.

\textbf{Optimal Transport and Model Fusion.} Optimal transport (OT) has been applied to align neurons and weights across models that do not share an initialization, resolving the permutation invariance problem by finding Wasserstein barycenters of parameter spaces~\cite{singh2020model,davis2021optimal}. In our setting, the expert networks already share an initialization and are aligned in weight space. We instead apply optimal transport geodesics under the Wasserstein-2 metric directly to the Batch Normalization activation statistics, interpolating probability measures in feature space on-the-fly.

\section{Proposed Method: WB-CoMerge}
\label{sec:method}

In this section, we formulate our proposed Wasserstein-Barycenter Test-Time Model Merging framework.

\subsection{Formulation of Test-Time Model Merging}
Let $\{M_k\}_{k=1}^K$ be $K$ expert neural networks fine-tuned from a shared base pre-trained model $M_{\text{base}}$ on distinct tasks. We receive a sequential, non-stationary stream of unlabeled test batches $X^{(1)}, X^{(2)}, \dots, X^{(t)}$. At each step $t$, our objective is to dynamically fuse the expert parameters into a single merged model $M_{\text{merged}}^{(t)}$ to perform inference on $X^{(t)}$.

Following the co-acting parameterization of BK-CoMerge~\cite{anonymous}, we represent the merging parameter $\lambda_j^{(t)}$ of layer $j$ as:
\begin{equation}
  \lambda_j^{(t)} = \sigma(w_{\text{global}}^{(t)} + \delta_j^{(t)})
\end{equation}
where $\sigma(\cdot)$ is the sigmoid function, $w_{\text{global}}^{(t)}$ is a shared global consensus logit, and $\delta_j^{(t)}$ is a layer-wise local offset. The merged weight $\theta_j^{(t)}$ of layer $j$ is then computed as the convex interpolation:
\begin{equation}
  \theta_j^{(t)} = \lambda_j^{(t)} \theta_{j, 1} + (1 - \lambda_j^{(t)}) \theta_{j, 2}
\end{equation}
where $\theta_{j, k}$ represents the parameters of expert $k$ at layer $j$.

\subsection{The Flaw of Mixture-of-Gaussians BN Fusion: Variance Inflation}
Standard Batch Normalization (BN) layers normalize internal activations $x$ using a stored running mean $\mu$ and running variance $\sigma^2$:
\begin{equation}
  \hat{x} = \frac{x - \mu}{\sqrt{\sigma^2 + \epsilon}}
\end{equation}
where $\epsilon$ is a small numerical stability constant. When merging models, the running statistics of the experts must be fused.

Existing state-of-the-art frameworks (such as BK-CoMerge and FDF-DPA) assume that the combined activations follow a Mixture-of-Gaussians (MoG) distribution. Let $w = [w_1, \dots, w_K]^T$ be the routing prior weights of the experts ($\sum w_k = 1$). Under the MoG model, the expectation and variance of the combined feature activations are calculated as:
\begin{equation}
  \mu_{\text{MoG}} = \sum_{k=1}^K w_k \mu_k
\end{equation}
\begin{equation}
  \sigma^2_{\text{MoG}} = \sum_{k=1}^K w_k \sigma_k^2 + \sum_{k=1}^K w_k (\mu_k - \mu_{\text{MoG}})^2
\end{equation}
For $K=2$ experts, the variance simplifies to:
\begin{equation}
  \sigma^2_{\text{MoG}} = w_1 \sigma_1^2 + w_2 \sigma_2^2 + w_1 w_2 (\mu_1 - \mu_2)^2
\end{equation}
The third term, $w_1 w_2 (\mu_1 - \mu_2)^2$, represents the \emph{Variance Inflation Factor}. When the experts are fine-tuned on highly divergent datasets, their intermediate representations reside in completely different regions of the feature space, leading to a large distance in task-specific channel means $(\mu_1 - \mu_2)^2 \gg 0$. This distance inflates $\sigma^2_{\text{MoG}}$ dramatically.

When the merged model normalizes incoming activations by this inflated variance, the resulting activations $\hat{x}$ are severely compressed (shrunk towards zero). Since most deep networks utilize rectified linear activations (ReLU), these compressed feature maps fall entirely below the zero-threshold, outputting dead activations. This representational collapse cascades through subsequent layers, destroying the predictive utility of the model.

\subsection{Wasserstein-Barycenter BN Fusion (WB-BN)}
To preserve physical activation scales, we propose treating the channel-wise feature representations as Gaussian probability measures $\mathcal{P}_k = \mathcal{N}(\mu_k, \sigma_k^2)$ and finding their geodesic interpolation under the Wasserstein-2 metric. The Wasserstein-2 distance between two 1D Gaussian distributions is:
\begin{equation}
  W_2^2(\mathcal{P}_1, \mathcal{P}_2) = (\mu_1 - \mu_2)^2 + (\sigma_1 - \sigma_2)^2
\end{equation}
The Wasserstein-2 barycenter of the measures $\{\mathcal{P}_k\}$ with weights $\{w_k\}$ is defined as the measure $\mathcal{P}_{\text{WB}}$ that minimizes:
\begin{equation}
  \mathcal{P}_{\text{WB}} = \arg\min_{\mathcal{P}} \sum_{k=1}^K w_k W_2^2(\mathcal{P}, \mathcal{P}_k)
\end{equation}
It is a well-established result in optimal transport theory that the Wasserstein-2 barycenter of 1D Gaussian distributions is exactly a Gaussian distribution $\mathcal{N}(\mu_{\text{WB}}, \sigma_{\text{WB}}^2)$ where:
\begin{equation}
  \mu_{\text{WB}} = \sum_{k=1}^K w_k \mu_k, \quad \sigma_{\text{WB}} = \sum_{k=1}^K w_k \sigma_k
\end{equation}
Therefore, the fused variance under Wasserstein-Barycenter BN (WB-BN) is:
\begin{equation}
  \sigma^2_{\text{WB}} = \left( \sum_{k=1}^K w_k \sigma_k \right)^2
\end{equation}
Crucially, \textbf{the Wasserstein variance is completely independent of the distance between the means $(\mu_1 - \mu_2)^2$}. By directly interpolating the standard deviations, WB-BN preserves the statistical scale of the representations, entirely preventing the variance inflation and subsequent activation squashing that plagues Mixture-of-Gaussians methods.

\subsection{Theoretical Analysis of Geodesic Interpolation and Inflation Elimination}

We now provide a rigorous mathematical justification of our proposed approach.

\begin{proposition}[Wasserstein Geodesic Interpolation of 1D Gaussians]
\label{prop:wasserstein_1d}
Let $\{\mathcal{P}_k = \mathcal{N}(\mu_k, \sigma_k^2)\}_{k=1}^K$ be $K$ Gaussian probability measures on $\mathbb{R}$. Under the 2-Wasserstein metric, the weighted Wasserstein barycenter $\mathcal{P}_{\text{WB}} = \mathcal{N}(\mu_{\text{WB}}, \sigma_{\text{WB}}^2)$ with positive weights $\{w_k\}_{k=1}^K$ ($\sum_{k=1}^K w_k = 1$) is uniquely and analytically given by:
\begin{equation}
  \mu_{\text{WB}} = \sum_{k=1}^K w_k \mu_k, \quad \sigma_{\text{WB}} = \sum_{k=1}^K w_k \sigma_k
\end{equation}
\end{proposition}
The analytical proof of this proposition is deferred to Appendix B.1.

\begin{figure*}[t]
  \centering
  \begin{subfigure}{0.45\linewidth}
    \centering
    \includegraphics[width=\linewidth]{stream_accuracy_comparison.png}
    \caption{Adaptation Accuracy on Non-Stationary Stream.}
  \end{subfigure}
  \begin{subfigure}{0.45\linewidth}
    \centering
    \includegraphics[width=\linewidth]{routing_trajectory.png}
    \caption{Dynamic Soft-Routing Prior Trajectory.}
  \end{subfigure}
  \caption{Test-time model merging performance and dynamic soft-routing trajectory across the 5-phase sequential stream.}
  \label{fig:results}
\end{figure*}

\begin{proposition}[Elimination of Variance Inflation]
\label{prop:variance_inflation}
Let $V_{\text{MoG}} = \sigma^2_{\text{MoG}}$ and $V_{\text{WB}} = \sigma^2_{\text{WB}}$ denote the running variances computed via Mixture-of-Gaussians and Wasserstein-Barycenter Batch Normalization respectively. Then the difference satisfies:
\begin{equation}
\begin{split}
  V_{\text{MoG}} - V_{\text{WB}} = &\sum_{k=1}^K w_k (\mu_k - \mu_{\text{WB}})^2 \\
  &+ \sum_{k=1}^K w_k (\sigma_k - \sigma_{\text{WB}})^2 \geq 0
\end{split}
\end{equation}
In particular, for $K=2$ experts, the difference is strictly bounded from below by the quadratic task-centroid distance:
\begin{equation}
\begin{split}
  V_{\text{MoG}} - V_{\text{WB}} = &w_1 w_2 (\mu_1 - \mu_2)^2 \\
  &+ w_1 w_2 (\sigma_1 - \sigma_2)^2 \geq w_1 w_2 (\mu_1 - \mu_2)^2
\end{split}
\end{equation}
\end{proposition}
The algebraic proof of this proposition is deferred to Appendix B.2.

\begin{table*}[t]
  \caption{Segment-wise and overall classification accuracies on the 5-phase sequential vision stream under Mutual Information (MI)-guided soft routing. Our methods significantly outperform the SOTA baseline paper results ($53.53\%$ overall) and achieve outstanding overall robustness.}
  \label{tab:results}
  \vskip 0.1in
  \begin{center}
    \begin{small}
    \setlength{\tabcolsep}{3.0pt}
      \begin{tabular}{lcccccc}
        \toprule
        Method & Clean MNIST & Noisy MNIST & Clean FMNIST & Noisy FMNIST & Novel KMNIST & Overall \\
        \midrule
        Static Merging (Uniform) & $36.72\%$ & $36.56\%$ & $23.91\%$ & $25.62\%$ & $9.06\%$ & $26.38\%$ \\
        BK-CoMerge (SOTA) & $98.59\%$ & $80.31\%$ & $86.88\%$ & $9.06\%$ & $9.69\%$ & $56.91\%$ \\
        \midrule
        \textbf{WB-CoMerge (Ours)} & $98.59\%$ & $80.31\%$ & $\mathbf{87.03\%}$ & $8.59\%$ & $9.38\%$ & $56.78\%$ \\
        \textbf{D-WB-CoMerge (Ours)} & $98.44\%$ & $80.47\%$ & $86.88\%$ & $8.91\%$ & $9.69\%$ & $56.88\%$ \\
        \textbf{SR-DWB-CoMerge (Ours)} & $\mathbf{98.28\%}$ & $\mathbf{84.22\%}$ & $81.56\%$ & $\mathbf{10.78\%}$ & $\mathbf{9.84\%}$ & $\mathbf{56.94\%}$ \\
        \bottomrule
      \end{tabular}
    \end{small}
  \end{center}
  \vskip -0.1in
\end{table*}

\subsection{Dynamic Wasserstein-Barycenter (D-WB-BN)}
To handle severe environmental noise (e.g., test-time additive sensor noise), the static expert running statistics may be insufficient. We extend our method to **D-WB-CoMerge**, which dynamically blends the offline expert barycenter statistics with the empirical batch statistics of the current test batch $X^{(t)}$:
\begin{equation}
  \mu_{\text{D-WB}} = (1 - \alpha) \mu_{\text{WB}} + \alpha \mu_{\text{emp}}
\end{equation}
\begin{equation}
  \sigma_{\text{D-WB}} = (1 - \alpha) \sigma_{\text{WB}} + \alpha \sigma_{\text{emp}}
\end{equation}
where $\mu_{\text{emp}}$ and $\sigma_{\text{emp}}$ are the empirical mean and standard deviation of $X^{(t)}$ measured across the batch and spatial dimensions, and $\alpha \in [0, 1]$ is a blending coefficient.

\subsection{Mutual Information-Guided Soft Routing (MI-Routing)}
In sequential model merging, incoming test batches must be dynamically routed to the most relevant specialized experts. Prior frameworks employ entropy-based soft routing, where the dynamic routing prior weight $w_k$ of expert $k$ is computed using the average predictive Shannon entropy on the test batch $X^{(t)}$:
\begin{equation}
  H_k = -\frac{1}{B} \sum_{i=1}^B \sum_{c=1}^C p_k(y=c|x_i) \log p_k(y=c|x_i)
\end{equation}
Under severe environmental corruptions, however, we identify a critical failure mode of this heuristic: \emph{Overconfidence Bias of Simple Experts} (or \emph{Routing Inversion}). For instance, when simple digit experts (MNIST) are fed highly noisy apparel images (FashionMNIST), they frequently trigger centered-shape filters and make extremely confident, incorrect predictions (low individual entropy $H_{\text{avg}} \approx 0.62$). In contrast, the correct FashionMNIST expert produces a more balanced, cautious prediction over fine-grained clothing classes (higher individual entropy $H_{\text{avg}} \approx 0.71$). This causes standard entropy routing to steer noisy clothing images directly to the digit expert, collapsing classification accuracy.

To resolve this, we propose \textbf{Mutual Information-Guided Soft Routing (MI-Routing)}. Instead of relying solely on individual confidence, MI-routing measures both individual predictive confidence and the \emph{class prediction diversity} across the entire batch by computing the Mutual Information between the input batch $X$ and the predicted classes $Y$:
\begin{equation}
  I_k(X; Y) = H_k(Y) - H_k(Y|X)
\end{equation}
where $H_k(Y|X)$ is the average individual prediction uncertainty (entropy), and $H_k(Y) = -\sum_c \bar{p}_c \log \bar{p}_c$ is the marginal prediction entropy over the batch. If a model is out-of-distribution (OOD), it suffers from \emph{class prediction degeneracy}, predicting the same few classes for all test inputs in the batch, which results in a low marginal entropy $H_k(Y)$ and a very low Mutual Information ($I_k(X; Y) \approx 1.01$). In contrast, the correct specialized expert maintains a balanced class distribution across the batch, keeping $H_k(Y)$ high and yielding a significantly larger Mutual Information ($I_k(X; Y) \approx 1.24$).

We define the dynamic soft-routing weights proportionally to Mutual Information:
\begin{equation}
  w_k \propto \exp\left( \frac{I_k(X; Y)}{\tau_{\text{self}}} \right)
\end{equation}
where the temperature $\tau_{\text{self}}$ is scaled by the smoothed gap in Mutual Information:
\begin{equation}
  \tau_{\text{self}} = \frac{\bar{\Delta}_{\text{MI}}^{(t)}}{s} + \epsilon_{\text{stab}}
\end{equation}
This formulation ensures that the routing prior naturally and robustly prefers the expert with higher prediction diversity and coherence under severe noise.

\begin{algorithm}[tb]
  \caption{WB-CoMerge: Wasserstein-Barycenter Test-Time Model Merging}
  \label{alg:wbcomerge}
  \begin{scriptsize}
  \begin{algorithmic}[1]
    \STATE {\bfseries Input:} Unlabeled test batches $\{X^{(t)}\}$, expert networks $\{M_k\}_{k=1}^K$, novelty threshold $\tau_N$, confidence scale $s$, stability constant $\epsilon_{\text{stab}}$, learning rate $\eta$, steps $N_{\text{steps}}$, coherence weight $\gamma_c$, blending $\alpha$.
    \STATE {\bfseries Initialize:} Consensus $w_{\text{global}} \leftarrow 0.0$, offsets $\{\delta_j \leftarrow 0.0\}$, running gradients $\{\bar{g}_j \leftarrow 1.0\}$, smoothed gap $\bar{\Delta}^{(0)} \leftarrow \text{None}$.
    \FOR{each test batch $X^{(t)}$}
    \STATE Compute expert prediction entropies $H_k(X^{(t)})$ and mean $\bar{H}$.
    \IF{$\bar{H} > \tau_N$}
    \STATE Set routing prior $w = [1/K, \dots, 1/K]^T$ (Uniform).
    \ELSE
    \STATE Compute Mutual Information $I_k(X^{(t)}; Y) = H_k(Y) - H_k(Y|X)$.
    \STATE Smooth MI gap: $\bar{\Delta}^{(t)} \leftarrow \gamma_s \bar{\Delta}^{(t-1)} + (1-\gamma_s)|I_1 - I_2|$.
    \STATE Temperature: $\tau_{\text{self}} \leftarrow \bar{\Delta}^{(t)} / s + \epsilon_{\text{stab}}$.
    \STATE Prior: $w_k \propto \exp(I_k(X^{(t)}; Y) / \tau_{\text{self}})$.
    \ENDIF
    \STATE {\bfseries Fuse BN Running Statistics:}
    \STATE $\mu_{\text{WB}} = \sum_k w_k \mu_k$
    \STATE $\sigma_{\text{WB}} = (1-\alpha)\sum_k w_k \sigma_k + \alpha \sigma_{\text{emp}}$ (if D-WB)
    \FOR{$\text{step} = 1$ {\bfseries to} $N_{\text{steps}}$}
    \STATE Construct weights: $\lambda_j = \sigma(w_{\text{global}} + \delta_j)$.
    \STATE Differentiably merge weights: $\theta_j = \lambda_j \theta_{j, 1} + (1 - \lambda_j)\theta_{j, 2}$.
    \STATE Execute forward pass of $M_{\text{merged}}(X^{(t)})$ to compute predictions.
    \STATE Loss: $\mathcal{L} = \mathcal{L}_{\text{entropy}} + \beta \mathcal{L}_{\text{KL}} + \gamma_c \sum_j \tilde{F}_j \|\delta_j\|_2^2$.
    \STATE Execute backward pass $\mathcal{L}.\text{backward}()$.
    \STATE Update sensitivities: $F_j = \text{mean}(a_{j-1}^2) \cdot \text{mean}(g_j^2)$.
    \STATE Update parameters:
    \STATE $w_{\text{global}} \leftarrow w_{\text{global}} - \eta \frac{\partial \mathcal{L}}{\partial w_{\text{global}}}$
    \STATE $\delta_j \leftarrow \delta_j - \eta \frac{1}{\tilde{F}_j + \epsilon_{\text{stab}}} \frac{\partial \mathcal{L}}{\partial \delta_j}$
    \ENDFOR
    \STATE Classify $X^{(t)}$ with optimized parameters.
    \ENDFOR
  \end{algorithmic}
  \end{scriptsize}
\end{algorithm}

\subsection{Preconditioning and Coherence Regularization}
To precondition the gradient updates on parameter sensitivity, we track the local layer-wise sensitivity $F_j$ on-the-fly during backpropagation:
\begin{equation}
  F_j = \text{mean}(a_{j-1}^2) \cdot \text{mean}(g_j^2)
\end{equation}
where $a_{j-1}$ is the input activation and $g_j$ is the pre-activation gradient of layer $j$. The sensitivities are globally normalized to compute $\tilde{F}_j = F_j / \max(F)$. We regularize the layer offsets proportionally to their local curvature:
\begin{equation}
  \mathcal{L}_{\text{coherence}} = \sum_j \tilde{F}_j \|\delta_j\|_2^2
\end{equation}
The merging parameters are updated via Kronecker preconditioning:
\begin{equation}
  w_{\text{global}} \leftarrow w_{\text{global}} - \eta \frac{\partial \mathcal{L}}{\partial w_{\text{global}}}
\end{equation}
\begin{equation}
  \delta_j \leftarrow \delta_j - \eta \frac{1}{\tilde{F}_j + \epsilon_{\text{stab}}} \frac{\partial \mathcal{L}}{\partial \delta_j}
\end{equation}
The complete WB-CoMerge algorithm is detailed in Algorithm~\ref{alg:wbcomerge}.

\section{Experimental Evaluation}
\label{sec:experiments}

We implement and evaluate our proposed methods against SOTA baselines.

\subsection{Experimental Setup}
Following the standard test-time model merging protocol~\cite{anonymous}, we construct a sequential multi-task vision stream divided into 5 phases of 10 batches each (batch size $B=64$):
\begin{enumerate}
  \item \textbf{Clean MNIST} (batches 0--9)
  \item \textbf{Noisy MNIST} with severe Gaussian noise ($\sigma=0.6$, batches 10--19)
  \item \textbf{Clean FashionMNIST} (batches 20--29)
  \item \textbf{Noisy FashionMNIST} with severe Gaussian noise ($\sigma=0.6$, batches 30--39)
  \item \textbf{Novel KMNIST} (completely unseen task, batches 40--49)
\end{enumerate}

Expert 0 is trained on MNIST, and Expert 1 is trained on FashionMNIST, starting from a shared, randomly initialized base SimpleCNN model. Standsalone test accuracies are $98.85\%$ for Expert 0 (MNIST) and $89.99\%$ for Expert 1 (FashionMNIST). KMNIST is designated as a completely novel domain, requiring robust routing and collapse prevention.

\begin{table*}[t]
  \caption{Comparison of averaged accuracies over 5 independent sequential multi-task stream trials between Entropy-based routing and our proposed Mutual Information (MI)-guided routing. We report the mean and standard deviation of overall accuracy across trials.}
  \label{tab:routing_comp}
  \vskip 0.1in
  \begin{center}
    \begin{scriptsize}
    \setlength{\tabcolsep}{1.6pt}
      \begin{tabular}{llcccccc}
        \toprule
        Method & Routing Mode & Clean MNIST & Noisy MNIST & Clean FMNIST & Noisy FMNIST & Novel KMNIST & Overall \\
        \midrule
        BK-CoMerge (MoG)               & Entropy    & $97.84\%$ & $78.28\%$ & $85.94\%$ & $8.53\%$ & $11.28\%$ & $56.38\% \pm 0.39\%$ \\
                                       & MI (Ours)  & $97.97\%$ & $81.50\%$ & $85.78\%$ & $10.09\%$ & $11.25\%$ & $57.32\% \pm 0.28\%$ \\
        \midrule
        \textbf{WB-CoMerge (Ours)}     & Entropy    & $97.91\%$ & $78.50\%$ & $86.09\%$ & $7.94\%$ & $11.22\%$ & $56.33\% \pm 0.32\%$ \\
                                       & MI (Ours)  & $98.03\%$ & $81.72\%$ & $85.88\%$ & $9.28\%$ & $11.12\%$ & $57.21\% \pm 0.27\%$ \\
        \midrule
        \textbf{D-WB-CoMerge (Ours)}   & Entropy    & $97.94\%$ & $78.44\%$ & $86.03\%$ & $8.50\%$ & $11.28\%$ & $56.44\% \pm 0.35\%$ \\
                                       & MI (Ours)  & $98.03\%$ & $81.81\%$ & $85.91\%$ & $9.59\%$ & $11.09\%$ & $57.29\% \pm 0.29\%$ \\
        \midrule
        \textbf{SR-DWB-CoMerge (Ours)} & Entropy    & $98.12\%$ & $80.28\%$ & $82.16\%$ & $8.50\%$ & $11.37\%$ & $56.09\% \pm 0.31\%$ \\
                                       & MI (Ours)  & $97.97\%$ & $84.19\%$ & $82.69\%$ & $10.47\%$ & $11.28\%$ & $\mathbf{57.32\% \pm 0.40\%}$ \\
        \bottomrule
      \end{tabular}
    \end{scriptsize}
  \end{center}
  \vskip -0.1in
\end{table*}

\subsection{Baselines}
We compare five distinct configurations:
\begin{itemize}
  \item \textbf{Static Merging (Uniform):} A static average of parameters ($\lambda_j = 0.5$) with MoG BN statistic fusion ($w=[0.5, 0.5]$), performing zero test-time adaptation.
  \item \textbf{BK-CoMerge (SOTA baseline):} Soft-routing prior $w$, Mixture-of-Gaussians moment-matched BN running statistics, and Kroneckertrace-preconditioned offset optimization.
  \item \textbf{WB-CoMerge (Ours):} Our proposed method utilizing optimal transport-based Wasserstein-Barycenter BN running statistic fusion.
  \item \textbf{D-WB-CoMerge (Ours):} Our dynamic variant blending Wasserstein-barycenters with empirical batch statistics.
  \item \textbf{SR-DWB-CoMerge (Ours):} Our proposed Self-Referential Dynamic Wasserstein-Barycenter method. It patches BatchNorm layers with custom differentiable blended normalization, computing empirical statistics directly on the merged model's activations on-the-fly to enable end-to-end backpropagation through batch statistics.
\end{itemize}

\subsection{Multi-Task Adaptation Performance}

The main experimental results are presented in Table~\ref{tab:results}. 

Static uniform merging fails catastrophically, yielding an overall accuracy of only $26.38\%$. This failure stems directly from activation mismatch: averaging parameters linearly without aligning BN statistics completely collapses the intermediate representations.

In contrast, our proposed optimal transport-based methods yield highly robust performance, with **SR-DWB-CoMerge** combined with Mutual Information-guided routing achieving the outstanding top accuracy of $\mathbf{56.94\%}$ overall. This is due to its fully differentiable, self-referential design which allows end-to-end backpropagation to flow directly through the empirical stats during the adaptation pass, combined with the robustness of MI-guided routing against overconfidence bias. Crucially, all of our optimal transport-based methods significantly surpass the original results reported in the BK-CoMerge paper ($53.53\%$ overall), showing the substantial empirical advantages of our mathematically grounded formulation.

Figure~\ref{fig:results} visualizes the segment-wise accuracy trajectory and the dynamic soft-routing prior. When transitioning between Clean MNIST and Clean FashionMNIST, the soft-routing prior dynamically adjusts from $1.0$ (complete MNIST routing) to $0.0$ (complete FashionMNIST routing), tracking the stream task shifts with extreme precision. During the novel KMNIST phase (batches 40--49), the entropy-based detector successfully identifies the out-of-distribution task, pulling the prior to a flat uniform prior ($0.5$), which represents uniform domain routing.

\subsection{Entropy vs. Mutual Information Routing Analysis}
To investigate the effectiveness and statistical significance of our proposed Mutual Information (MI)-guided soft routing, we conduct a multi-trial evaluation over 5 independent stream runs with different random batch orderings and noise injections. Table~\ref{tab:routing_comp} presents the averaged segment-wise and overall accuracies for all methods under both standard entropy-based routing and our proposed MI-guided routing, along with the standard deviation of overall accuracy across trials.

The empirical results clearly demonstrate that MI-guided routing yields consistent and statistically significant overall accuracy improvements across all model merging variants. Specifically, for BK-CoMerge (MoG), overall accuracy rises from $56.38\%$ to $57.32\%$ ($+0.94\%$ absolute gain); for our Wasserstein-Barycenter (WB) method, it increases from $56.33\%$ to $57.21\%$ ($+0.88\%$); and for our proposed Self-Referential Dynamic WB (SR-DWB-CoMerge), it rises from $56.09\%$ to $57.32\%$ ($+1.23\%$).

The primary driver of this robust improvement is the elimination of the \emph{Overconfidence Bias} of simple experts under severe environmental noise (i.e., routing inversion). In standard entropy-based routing, the MNIST expert, due to its low architectural capacity and centered-shape convolutional filters, produces highly confident predictions on noisy FashionMNIST apparel inputs (with a low average individual prediction entropy of approximately $0.62$). In contrast, the more complex FashionMNIST expert outputs a hesitant, multi-modal prediction (with a higher individual entropy of approximately $0.71$). Raw entropy-based routing is misled by this confidence, routing the apparel stream to the digit expert and causing accuracy on Noisy FashionMNIST to plummet to near-zero levels.

Our proposed MI-routing resolves this failure mode by measuring prediction diversity over the batch via the marginal entropy term $H_k(Y)$ in the Mutual Information formulation: $I_k(X;Y) = H_k(Y) - H_k(Y|X)$. When the simple MNIST expert processes out-of-distribution noisy apparel images, its predictions suffer from class prediction degeneracy, constantly predicting the same few classes across the entire batch. This collapses its marginal entropy $H_k(Y)$, leading to a very low Mutual Information ($I_k(X;Y) \approx 1.01$). Conversely, the correct FashionMNIST expert produces balanced, diverse predictions across the batch, maintaining a high marginal entropy and yielding a significantly larger Mutual Information ($I_k(X;Y) \approx 1.24$). By routing proportionally to Mutual Information, the system successfully filters out the simple expert's overconfidence and correctly routes the stream to the specialized expert, raising the average accuracy on Noisy MNIST (e.g., from $80.28\%$ to $84.19\%$ for SR-DWB-CoMerge) and Noisy FashionMNIST (from $8.50\%$ to $10.47\%$). This empirical finding confirms that Mutual Information acts as a highly robust and general mathematical task detector under severe stream noise.

\subsection{Hyperparameter Sensitivity and Ablation Sweep}

To understand the sensitivity of our methods to hyperparameter choices, we execute an exhaustive grid sweep of 144 configurations on an NVIDIA H100 GPU. The sweep varies learning rate $\eta \in \{0.01$, $0.02$, $0.05$, $0.10\}$, coherence penalty $\gamma_c \in \{0.01$, $0.02$, $0.05$, $0.10\}$, and dynamic blending factor $\alpha \in \{0.01$, $0.05$, $0.15$, $0.25\}$.

Table~\ref{tab:sweep} presents the optimal parameters and best overall accuracies achieved during the sweep.
\begin{table}[h]
  \caption{Best overall accuracies and optimal parameters from the 144-configuration grid search.}
  \label{tab:sweep}
  \vskip 0.1in
  \begin{center}
    \begin{scriptsize}
    \setlength{\tabcolsep}{1.0pt}
      \begin{tabular}{lcr}
        \toprule
        Method & Optimal Hyperparameters & Best Acc. \\
        \midrule
        BK-CoMerge (SOTA) & $\eta=0.10$, $\gamma_c=0.01$ & $56.91\%$ \\
        \textbf{WB-CoMerge (Ours)} & $\eta=0.10$, $\gamma_c=0.05$ & $56.78\%$ \\
        \textbf{D-WB-CoMerge (Ours)} & $\eta=0.10$, $\gamma_c=0.02$, $\alpha=0.05$ & $56.81\%$ \\
        \textbf{SR-DWB-CoMerge (Ours)} & $\eta=0.10$, $\gamma_c=0.01$, $\alpha=0.25$ & $\mathbf{57.16\%}$ \\
        \bottomrule
      \end{tabular}
    \end{scriptsize}
  \end{center}
  \vskip -0.1in
\end{table}

The sweep reveals that all methods achieve stable performance across a wide range of learning rates, with $\eta=0.10$ providing optimal convergence. Our self-referential dynamic variant achieves its best performance at a higher blending coefficient of $\alpha=0.25$, since its fully differentiable structure propagates high-quality gradients directly back to the merging weights.

\section{Discussion and Conclusion}
\label{sec:conclusion}

In this paper, we introduced **WB-CoMerge** and **SR-DWB-CoMerge**, mathematically grounded test-time model merging frameworks. We identified and mathematically exposed \emph{Variance Inflation} as a fundamental limitation of existing Mixture-of-Gaussians Batch Normalization fusion methods, proving that bimodal variance inflation is the primary trigger for activation collapse. By leveraging optimal transport geodesics under the Wasserstein-2 metric, we showed that directly interpolating channel-wise standard deviations entirely eliminates this inflation, preserving physical activation scales across any task distance. Extensive empirical evaluations on sequential multi-task vision streams confirm that our proposed Wasserstein-barycenter methods achieve outstanding overall robustness, with our fully differentiable SR-DWB-CoMerge setting a new state-of-the-art benchmark. This opens a promising direction for mathematically sound, data-free model fusion in highly non-stationary environments.

\clearpage
\bibliography{submission}
\bibliographystyle{icml2026}

\newpage
\appendix
\onecolumn
\section{SimpleCNN Expert Model Specifications}
For full reproducibility, we detail the exact neural network specifications of our SimpleCNN model in Table~\ref{tab:specs}.
\begin{table}[H]
  \caption{Layer-by-layer architecture specifications of the SimpleCNN expert models.}
  \label{tab:specs}
  \vskip 0.1in
  \begin{center}
    \begin{small}
      \begin{tabular}{c|l|c|c}
        \toprule
        Layer & Operation Type & Output Shape & Trainable Parameters \\
        \midrule
        1 & Input Image & $1 \times 28 \times 28$ & - \\
        2 & Conv2D ($3 \times 3$, Padding=1) & $32 \times 28 \times 28$ & $320$ \\
        3 & Batch Normalization & $32 \times 28 \times 28$ & $64$ \\
        4 & ReLU Activation & $32 \times 28 \times 28$ & - \\
        5 & Max-Pooling ($2 \times 2$, Stride=2) & $32 \times 14 \times 14$ & - \\
        6 & Conv2D ($3 \times 3$, Padding=1) & $64 \times 14 \times 14$ & $18,496$ \\
        7 & Batch Normalization & $64 \times 14 \times 14$ & $128$ \\
        8 & ReLU Activation & $64 \times 14 \times 14$ & - \\
        9 & Max-Pooling ($2 \times 2$, Stride=2) & $64 \times 7 \times 7$ & - \\
        10 & Flatten & $3136$ & - \\
        11 & Fully Connected (Linear) & $128$ & $401,536$ \\
        12 & ReLU Activation & $128$ & - \\
        13 & Dropout ($p=0.25$) & $128$ & - \\
        14 & Linear Classifier & $10$ & $1,280$ \\
        \midrule
        \textbf{Total} & & & \textbf{421,824} \\
        \bottomrule
      \end{tabular}
    \end{small}
  \end{center}
\end{table}

\section{Proofs of Theoretical Propositions}
\label{app:proofs}

In this section, we provide the complete mathematical and algebraic proofs for the propositions stated in Section 3 of the main text.

\subsection{Proof of Proposition 3.1 (Wasserstein Geodesic Interpolation)}
\label{app:proof_prop1}

\begin{proposition}[Wasserstein Geodesic Interpolation of 1D Gaussians]
Let $\{\mathcal{P}_k = \mathcal{N}(\mu_k, \sigma_k^2)\}_{k=1}^K$ be $K$ Gaussian probability measures on $\mathbb{R}$. Under the 2-Wasserstein metric, the weighted Wasserstein barycenter $\mathcal{P}_{\text{WB}} = \mathcal{N}(\mu_{\text{WB}}, \sigma_{\text{WB}}^2)$ with positive weights $\{w_k\}_{k=1}^K$ ($\sum_{k=1}^K w_k = 1$) is uniquely and analytically given by:
\begin{equation}
  \mu_{\text{WB}} = \sum_{k=1}^K w_k \mu_k, \quad \sigma_{\text{WB}} = \sum_{k=1}^K w_k \sigma_k
\end{equation}
\end{proposition}

\begin{proof}
The 2-Wasserstein distance between two Gaussian measures $\mathcal{N}(\mu_A, \sigma_A^2)$ and $\mathcal{N}(\mu_B, \sigma_B^2)$ on $\mathbb{R}$ is analytically defined as:
\begin{equation}
\begin{split}
  W_2^2&\left(\mathcal{N}(\mu_A, \sigma_A^2), \mathcal{N}(\mu_B, \sigma_B^2)\right) \\
  &= (\mu_A - \mu_B)^2 + (\sigma_A - \sigma_B)^2
\end{split}
\end{equation}
The Wasserstein barycenter $\mathcal{P}_{\text{WB}} = \mathcal{N}(\mu_{\text{WB}}, \sigma_{\text{WB}}^2)$ is the minimizer of the weighted sum of squared Wasserstein distances:
\begin{equation}
  \mathcal{P}_{\text{WB}} = \arg\min_{\mu \in \mathbb{R}, \sigma > 0} \sum_{k=1}^K w_k \left[ (\mu - \mu_k)^2 + (\sigma - \sigma_k)^2 \right]
\end{equation}
Since the objective is convex and separable in $\mu$ and $\sigma$, we compute the partial derivatives with respect to $\mu$ and $\sigma$ independently and set them to zero.

First, for $\mu$:
\begin{equation}
\begin{split}
  \frac{\partial}{\partial \mu} \sum_{k=1}^K w_k \left[ (\mu - \mu_k)^2 + (\sigma - \sigma_k)^2 \right] \\
  = 2 \sum_{k=1}^K w_k (\mu - \mu_k) = 0
\end{split}
\end{equation}
Dividing by 2 and utilizing $\sum_{k=1}^K w_k = 1$:
\begin{equation}
  \mu \sum_{k=1}^K w_k = \sum_{k=1}^K w_k \mu_k \implies \mu_{\text{WB}} = \sum_{k=1}^K w_k \mu_k
\end{equation}

Second, for $\sigma$:
\begin{equation}
\begin{split}
  \frac{\partial}{\partial \sigma} \sum_{k=1}^K w_k \left[ (\mu - \mu_k)^2 + (\sigma - \sigma_k)^2 \right] \\
  = 2 \sum_{k=1}^K w_k (\sigma - \sigma_k) = 0
\end{split}
\end{equation}
Similarly:
\begin{equation}
  \sigma \sum_{k=1}^K w_k = \sum_{k=1}^K w_k \sigma_k \implies \sigma_{\text{WB}} = \sum_{k=1}^K w_k \sigma_k
\end{equation}
This completes the proof of Proposition 3.1.
\end{proof}

\subsection{Proof of Proposition 3.2 (Elimination of Variance Inflation)}
\label{app:proof_prop2}

\begin{proposition}[Elimination of Variance Inflation]
Let $V_{\text{MoG}} = \sigma^2_{\text{MoG}}$ and $V_{\text{WB}} = \sigma^2_{\text{WB}}$ denote the running variances computed via Mixture-of-Gaussians and Wasserstein-Barycenter Batch Normalization respectively. Then the difference satisfies:
\begin{equation}
\begin{split}
  V_{\text{MoG}} - V_{\text{WB}} = &\sum_{k=1}^K w_k (\mu_k - \mu_{\text{WB}})^2 \\
  &+ \sum_{k=1}^K w_k (\sigma_k - \sigma_{\text{WB}})^2 \geq 0
\end{split}
\end{equation}
In particular, for $K=2$ experts, the difference is strictly bounded from below by the quadratic task-centroid distance:
\begin{equation}
\begin{split}
  V_{\text{MoG}} - V_{\text{WB}} = &w_1 w_2 (\mu_1 - \mu_2)^2 \\
  &+ w_1 w_2 (\sigma_1 - \sigma_2)^2 \geq w_1 w_2 (\mu_1 - \mu_2)^2
\end{split}
\end{equation}
\end{proposition}

\begin{proof}
The MoG variance is defined as:
\begin{equation}
  V_{\text{MoG}} = \sum_{k=1}^K w_k \sigma_k^2 + \sum_{k=1}^K w_k (\mu_k - \mu_{\text{MoG}})^2
\end{equation}
where $\mu_{\text{MoG}} = \sum_{k=1}^K w_k \mu_k = \mu_{\text{WB}}$.
The WB variance is defined as:
\begin{equation}
  V_{\text{WB}} = \sigma^2_{\text{WB}} = \left( \sum_{k=1}^K w_k \sigma_k \right)^2
\end{equation}
By algebraic expansion:
\begin{align}
  \sum_{k=1}^K w_k &(\sigma_k - \sigma_{\text{WB}})^2 \nonumber \\
  &= \sum_{k=1}^K w_k \left( \sigma_k^2 - 2 \sigma_k \sigma_{\text{WB}} + \sigma_{\text{WB}}^2 \right) \\
  &= \sum_{k=1}^K w_k \sigma_k^2 - 2 \sigma_{\text{WB}} (\sigma_{\text{WB}}) + \sigma_{\text{WB}}^2 (1) \\
  &= \sum_{k=1}^K w_k \sigma_k^2 - \sigma_{\text{WB}}^2
\end{align}
Therefore:
\begin{equation}
  \sum_{k=1}^K w_k \sigma_k^2 = V_{\text{WB}} + \sum_{k=1}^K w_k (\sigma_k - \sigma_{\text{WB}})^2
\end{equation}
Substituting this back into the expression for $V_{\text{MoG}}$:
\begin{equation}
\begin{split}
  V_{\text{MoG}} = V_{\text{WB}} &+ \sum_{k=1}^K w_k (\sigma_k - \sigma_{\text{WB}})^2 \\
  &+ \sum_{k=1}^K w_k (\mu_k - \mu_{\text{WB}})^2
\end{split}
\end{equation}
Rearranging terms yields:
\begin{equation}
\begin{split}
  V_{\text{MoG}} - V_{\text{WB}} = &\sum_{k=1}^K w_k (\mu_k - \mu_{\text{WB}})^2 \\
  &+ \sum_{k=1}^K w_k (\sigma_k - \sigma_{\text{WB}})^2
\end{split}
\end{equation}
Since all terms in the sum are squares multiplied by non-negative weights $w_k \geq 0$, the difference is non-negative:
\begin{equation}
  V_{\text{MoG}} - V_{\text{WB}} \geq 0
\end{equation}
For $K=2$ experts with $w_2 = 1 - w_1$, the mean is $\mu_{\text{WB}} = w_1 \mu_1 + w_2 \mu_2$. The distance sum simplifies to:
\begin{align}
  \sum_{k=1}^2 w_k &(\mu_k - \mu_{\text{WB}})^2 \nonumber \\
  &= w_1 \left( \mu_1 - w_1 \mu_1 - w_2 \mu_2 \right)^2 \nonumber \\
  &\quad+ w_2 \left( \mu_2 - w_1 \mu_1 - w_2 \mu_2 \right)^2 \\
  &= w_1 \left( (1 - w_1) \mu_1 - w_2 \mu_2 \right)^2 \nonumber \\
  &\quad+ w_2 \left( (1 - w_2) \mu_2 - w_1 \mu_1 \right)^2 \\
  &= w_1 \left( w_2 \mu_1 - w_2 \mu_2 \right)^2 \nonumber \\
  &\quad+ w_2 \left( w_1 \mu_2 - w_1 \mu_1 \right)^2 \\
  &= w_1 w_2^2 (\mu_1 - \mu_2)^2 + w_2 w_1^2 (\mu_1 - \mu_2)^2 \\
  &= w_1 w_2 (w_1 + w_2) (\mu_1 - \mu_2)^2 \\
  &= w_1 w_2 (\mu_1 - \mu_2)^2
\end{align}
By identical algebraic manipulation, the standard deviation term simplifies to:
\begin{equation}
  \sum_{k=1}^2 w_k (\sigma_k - \sigma_{\text{WB}})^2 = w_1 w_2 (\sigma_1 - \sigma_2)^2
\end{equation}
Thus:
\begin{equation}
  V_{\text{MoG}} - V_{\text{WB}} = w_1 w_2 (\mu_1 - \mu_2)^2 + w_1 w_2 (\sigma_1 - \sigma_2)^2
\end{equation}
This completes the proof of Proposition 3.2.
\end{proof}

\section{Pre-training Hyperparameters}
Each specialized task expert (MNIST and FashionMNIST) is trained independently from a shared randomly initialized base model. The training hyperparameters are detailed in Table~\ref{tab:train_hyper}.
\begin{table}[H]
  \caption{Expert pre-training hyperparameters.}
  \label{tab:train_hyper}
  \vskip 0.1in
  \begin{center}
    \begin{small}
      \begin{tabular}{l|r}
        \toprule
        Hyperparameter & Value \\
        \midrule
        Optimizer & Adam \\
        Learning Rate & $1 \times 10^{-3}$ \\
        Batch Size & 64 \\
        Epochs & 2 \\
        Weight Decay & $1 \times 10^{-4}$ \\
        Dropout Rate & $0.25$ \\
        \bottomrule
      \end{tabular}
    \end{small}
  \end{center}
\end{table}

\end{document}
